\section{Troubleshooting and known limitations}
\label{sec:trouble}

\subsection{Symptoms and what they mean}

\begin{longtable}{@{}p{5.4cm}p{9cm}@{}}
\toprule
Symptom & Cause and remedy \\
\midrule
\endhead
The first solve takes ten seconds & The scrambler tables are being built.  They are
cached for the life of the process; call \texttt{preload()} at start-up to pay the
cost once, on a thread of your own. \\
\texttt{error: not WCA move notation} & A token in the scramble is not a move.  The
message names the token and its index; \texttt{3Rw} and \texttt{1Rw} are invalid on
a $3\times3\times3$. \\
\texttt{error: the corner orientation sum must be divisible by three} & The state is
impossible: a single corner is twisted.  Re-check the facelets, or the tool that
produced them. \\
\texttt{error: each color must appear exactly nine times} & A facelet string with
nine of each color is the minimum requirement; a missing or duplicated sticker
usually means a transcription slip. \\
\texttt{error: the last layer does not match any case} & The last layer is not in a
state the tables know.  The pipeline normally catches this and finishes with the
two-phase solver; seeing the error directly means a table file is damaged. \\
\texttt{SearchLimitError} & A search was given too small a depth budget.  Raise
\texttt{SearchLimits.max\_total\_depth} (the default of 30 is generous) or let the
pair search run to its default depth of twelve. \\
\texttt{SolveTimeoutError} & A solve exceeded \texttt{SearchLimits.timeout\_ms} or
\texttt{SearchLimits.max\_nodes}.  Scramble generation gives each attempt a budget of
400,000 search nodes (about a second) and redraws a state that overruns; counting
nodes rather than time keeps a seed's scramble identical on every machine. \\
\texttt{CancelledError} & Someone called \texttt{CancelToken.cancel()}.  Expected
when an interface abandons a solve; reuse the token with \texttt{reset()}. \\
Solving feels slow & The F2L search is doing the work: the 41-case table only
covers configurations whose pieces are in the U layer or the slot.  Correct, but a
second or two per pair. \\
The net looks upside down & It is: the cross is at the bottom.  The
\texttt{Orientation} line says how the cube was turned to get there. \\
\bottomrule
\end{longtable}

\subsection{Known limitations}

\begin{enumerate}[leftmargin=1.4em]
  \item \textbf{F2L speed.}  The largest available improvement.  A complete table
  over the 576 pair configurations (or an explicit extraction step before the case
  table) would remove the search from the common path.
  \item \textbf{Last-layer naming.}  For a small number of states the stored
  algorithm does not finish the cube; the pipeline completes it with the two-phase
  solver instead of a named case.  The solution is still correct and verified, but
  the stage is reported as such rather than as \texttt{OLL (\textit{n})}.
  \item \textbf{Rotations inside algorithms.}  A few published algorithms include a
  rotation or a wide or slice move.  The stage re-orients the display afterwards so
  the cube stays cross-down; the solution remains valid, and the printed move list
  contains those rotations.
  \item \textbf{One puzzle.}  Only the $3\times3\times3$ is supported.  The cube
  model, the notation layer and the notation grammar are written for a cube of six
  faces and would need work for larger puzzles.
\end{enumerate}

\subsection{Verifying an installation}

\begin{lstlisting}[style=shell]
scripts/check.sh          # ruff format --check, ruff check, full pytest
uv run rubiks-cfop --scramble "R U R' U' F2 D" --no-color
\end{lstlisting}

The test suite covers the move tables (against Kociemba's published permutations),
the parser, the cubie invariants, both scramblers, the cross solver's optimality,
the F2L tables, exhaustive coverage of the last layer, the pipeline's stage order
and output format, and the command line itself.  A green run is the quickest proof
that the installation is sound.
